\chapter{R\^oles et Droits}
\label{chap:roles}

\section{Pr\'esentation des r\^oles}

L'application CGP met en \oe uvre un contr\^ole d'acc\`es bas\'e sur les r\^oles (\textit{Role-Based Access Control}). Chaque utilisateur se voit attribuer un ou plusieurs r\^oles via des \textbf{affectations}. Un r\^ole est toujours li\'e \`a une \textbf{structure} et \`a une \textbf{ann\'ee acad\'emique}.

\medskip
\begin{notebox}
  Un m\^eme utilisateur peut poss\'eder plusieurs r\^oles dans des structures diff\'erentes, voire dans la m\^eme structure. Les droits dont il dispose correspondent \`a l'union de tous ses r\^oles actifs.
\end{notebox}

\section{Catalogue des r\^oles}

Le tableau suivant pr\'esente l'ensemble des r\^oles disponibles dans l'application, class\'es par niveau d'acc\`es d\'ecroissant.

\medskip
\begin{longtable}{L{4.5cm}L{3cm}X}
  \toprule
  \rowcolor{cgplightblue}
  \textbf{Identifiant technique} & \textbf{Niveau} & \textbf{Description} \\
  \midrule
  \endfirsthead
  \toprule
  \rowcolor{cgplightblue}
  \textbf{Identifiant technique} & \textbf{Niveau} & \textbf{Description} \\
  \midrule
  \endhead
  \midrule
  \multicolumn{3}{r}{\small\textit{(suite page suivante)}}\\
  \endfoot
  \bottomrule
  \endlastfoot

  \texttt{services-centraux} & \cellcolor{cgplightred}Maximal
    & Acc\`es complet sans restriction. G\`ere les ann\'ees, les structures, les imports/exports, l'audit, et tous les param\'etrages. \\
  \rowcolor{cgplightgray}
  \texttt{administrateur} & Administratif
    & Administration du syst\`eme\,: gestion des utilisateurs, imports/exports, d\'el\'egations, organigrammes. \\
  \texttt{directeur-composante} & Direction composante
    & Gestion des utilisateurs et affectations dans sa composante. Import/export, organigrammes, d\'el\'egations dans son p\'erim\`etre. \\
  \rowcolor{cgplightgray}
  \texttt{directeur-administratif} & Direction admin.
    & Gestion administrative\,: utilisateurs, affectations, organigrammes, d\'el\'egations dans son p\'erim\`etre. \\
  \texttt{directeur-administratif-adjoint} & Direction admin. adj.
    & Similaire \`a \texttt{directeur-administratif} avec des droits l\'eg\`erement r\'eduits. \\
  \rowcolor{cgplightgray}
  \texttt{directeur-departement} & Direction d\'ept.
    & Gestion des affectations et d\'el\'egations dans son d\'epartement. G\'en\`ere des organigrammes. \\
  \texttt{vice-president-departement} & VP d\'ept.
    & Similaire au directeur de d\'epartement, avec acc\`es aux organigrammes et d\'el\'egations. \\
  \rowcolor{cgplightgray}
  \texttt{directeur-adjoint-licence} & Direction adj. licence
    & Gestion du p\'erim\`etre licence. Affectations et organigrammes. \\
  \texttt{responsable-service-pedagogique} & Resp. p\'edagogique
    & Gestion des affectations dans le service p\'edagogique. Organigrammes. \\
  \rowcolor{cgplightgray}
  \texttt{responsable-adjoint-service-pedagogique} & Resp. adj. p\'edagogique
    & Droits r\'eduits par rapport au responsable du service p\'edagogique. \\
  \texttt{directeur-mention} & Direction mention
    & Gestion de la mention. Affectations, organigrammes, d\'el\'egations. \\
  \rowcolor{cgplightgray}
  \texttt{directeur-specialite} & Direction sp\'ecialit\'e
    & Gestion de la sp\'ecialit\'e (mention). Affectations, organigrammes. \\
  \texttt{responsable-formation} & Resp. formation
    & Responsable d'une formation. Affectations li\'ees \`a sa formation, organigrammes. \\
  \rowcolor{cgplightgray}
  \texttt{responsable-annee} & Resp. ann\'ee
    & Responsable d'une ann\'ee d'\'etude. Affectations li\'ees, organigrammes. \\
  \texttt{directeur-etudes} & Directeur \'etudes
    & G\`ere les affectations li\'ees aux \'etudes. Organigrammes. \\
  \rowcolor{cgplightgray}
  \texttt{responsable-qualite} & Resp. qualit\'e
    & Acc\`es \`a l'annuaire et aux organigrammes dans son p\'erim\`etre. \\
  \texttt{responsable-international} & Resp. international
    & Acc\`es \`a l'annuaire et aux organigrammes li\'es aux programmes internationaux. \\
  \rowcolor{cgplightgray}
  \texttt{referent-commun} & R\'ef\'erent commun
    & Acc\`es g\'en\'eral \`a l'annuaire et aux organigrammes. \\
  \texttt{directeur-adjoint-ecole} & Directeur adj. \'ecole
    & Gestion du p\'erim\`etre \'ecole. Affectations et organigrammes. \\
  \rowcolor{cgplightgray}
  \texttt{secretariat-pedagogique} & Secr\'etariat
    & Consultation de l'annuaire. Cr\'eation et mise \`a jour des contacts de r\^ole. Signalements. \\
  \texttt{utilisateur-simple} & Utilisateur
    & Consultation de l'annuaire et des organigrammes. Signalements. \\
  \rowcolor{cgplightgray}
  \texttt{lecture-seule} & Lecture seule
    & Consultation uniquement. Aucune modification. Signalements. \\

\end{longtable}

\section{Regroupement fonctionnel des r\^oles}

Pour faciliter la lecture de ce manuel, les r\^oles sont regroup\'es en \textbf{profils fonctionnels} correspondant aux chapitres d\'edi\'es\,:

\medskip
\begin{tabularx}{\textwidth}{L{3.5cm}XC{2.5cm}}
  \toprule
  \rowcolor{cgplightblue}
  \textbf{Profil fonctionnel} & \textbf{R\^oles inclus} & \textbf{Chapitre} \\
  \midrule
  Services Centraux
    & \texttt{services-centraux}
    & Chap.~\ref{chap:sc} \\
  \rowcolor{cgplightgray}
  Direction composante / admin.
    & \texttt{directeur-composante}, \texttt{administrateur}, \texttt{directeur-administratif}, \texttt{directeur-administratif-adjoint}
    & Chap.~\ref{chap:dircomp} \\
  Directions et responsables
    & \texttt{directeur-departement}, \texttt{vice-president-departement}, \texttt{directeur-adjoint-licence}, \texttt{responsable-service-pedagogique}, \texttt{responsable-adjoint-service-pedagogique}, \texttt{directeur-mention}, \texttt{directeur-specialite}, \texttt{responsable-formation}, \texttt{responsable-annee}, \texttt{directeur-etudes}, \texttt{responsable-qualite}, \texttt{responsable-international}, \texttt{referent-commun}, \texttt{directeur-adjoint-ecole}
    & Chap.~\ref{chap:directions} \\
  \rowcolor{cgplightgray}
  Secr\'etariat p\'edagogique
    & \texttt{secretariat-pedagogique}
    & Chap.~\ref{chap:secretariat} \\
  Utilisateur standard
    & \texttt{utilisateur-simple}, \texttt{lecture-seule}
    & Chap.~\ref{chap:utilisateur} \\
  \bottomrule
\end{tabularx}

\section{Matrice des droits}

Le tableau suivant synth\'etise les droits par fonctionnalit\'e et par profil. \\
L\'egende\,: \textbf{\textcolor{cgpgreen}{\checkmark}} = autoris\'e, \textbf{\textcolor{cgpred}{$\times$}} = refus\'e, \textbf{\textcolor{cgporange}{\textasciitilde}} = partiel (p\'erim\`etre limit\'e).

\medskip
\begin{longtable}{L{4.5cm}C{1.6cm}C{1.9cm}C{2cm}C{1.8cm}C{1.5cm}}
  \toprule
  \rowcolor{cgplightblue}
  \textbf{Fonctionnalit\'e}
    & \textbf{SC}
    & \textbf{Dir. comp.}
    & \textbf{Directions}
    & \textbf{Secr\'et.}
    & \textbf{Util.} \\
  \midrule
  \endfirsthead
  \toprule
  \rowcolor{cgplightblue}
  \textbf{Fonctionnalit\'e}
    & \textbf{SC}
    & \textbf{Dir. comp.}
    & \textbf{Directions}
    & \textbf{Secr\'et.}
    & \textbf{Util.} \\
  \midrule
  \endhead
  \bottomrule
  \endlastfoot

  Consulter l'annuaire
    & {\color{cgpgreen}\checkmark}
    & {\color{cgpgreen}\checkmark}
    & {\color{cgpgreen}\checkmark}
    & {\color{cgpgreen}\checkmark}
    & {\color{cgpgreen}\checkmark} \\
  \rowcolor{cgplightgray}
  Voir les organigrammes
    & {\color{cgpgreen}\checkmark}
    & {\color{cgpgreen}\checkmark}
    & {\color{cgpgreen}\checkmark}
    & {\color{cgpgreen}\checkmark}
    & {\color{cgpgreen}\checkmark} \\
  G\'en\'erer des organigrammes
    & {\color{cgpgreen}\checkmark}
    & {\color{cgporange}\textasciitilde}
    & {\color{cgporange}\textasciitilde}
    & {\color{cgpred}$\times$}
    & {\color{cgpred}$\times$} \\
  \rowcolor{cgplightgray}
  Geler/d\'egeler un organigramme
    & {\color{cgpgreen}\checkmark}
    & {\color{cgpred}$\times$}
    & {\color{cgpred}$\times$}
    & {\color{cgpred}$\times$}
    & {\color{cgpred}$\times$} \\
  Cr\'eer des utilisateurs
    & {\color{cgpgreen}\checkmark}
    & {\color{cgporange}\textasciitilde}
    & {\color{cgpred}$\times$}
    & {\color{cgpred}$\times$}
    & {\color{cgpred}$\times$} \\
  \rowcolor{cgplightgray}
  Modifier des utilisateurs
    & {\color{cgpgreen}\checkmark}
    & {\color{cgporange}\textasciitilde}
    & {\color{cgpred}$\times$}
    & {\color{cgpred}$\times$}
    & {\color{cgpred}$\times$} \\
  Supprimer des utilisateurs
    & {\color{cgpgreen}\checkmark}
    & {\color{cgporange}\textasciitilde}
    & {\color{cgpred}$\times$}
    & {\color{cgpred}$\times$}
    & {\color{cgpred}$\times$} \\
  \rowcolor{cgplightgray}
  G\'erer les affectations
    & {\color{cgpgreen}\checkmark}
    & {\color{cgporange}\textasciitilde}
    & {\color{cgporange}\textasciitilde}
    & {\color{cgpred}$\times$}
    & {\color{cgpred}$\times$} \\
  G\'erer les contacts de r\^ole
    & {\color{cgpgreen}\checkmark}
    & {\color{cgpgreen}\checkmark}
    & {\color{cgpgreen}\checkmark}
    & {\color{cgpgreen}\checkmark}
    & {\color{cgpred}$\times$} \\
  \rowcolor{cgplightgray}
  Modifier les structures
    & {\color{cgpgreen}\checkmark}
    & {\color{cgpred}$\times$}
    & {\color{cgpred}$\times$}
    & {\color{cgpred}$\times$}
    & {\color{cgpred}$\times$} \\
  Importer des donn\'ees
    & {\color{cgpgreen}\checkmark}
    & {\color{cgporange}\textasciitilde}
    & {\color{cgpred}$\times$}
    & {\color{cgpred}$\times$}
    & {\color{cgpred}$\times$} \\
  \rowcolor{cgplightgray}
  Exporter des donn\'ees
    & {\color{cgpgreen}\checkmark}
    & {\color{cgporange}\textasciitilde}
    & {\color{cgpred}$\times$}
    & {\color{cgpred}$\times$}
    & {\color{cgpred}$\times$} \\
  G\'erer les ann\'ees acad\'emiques
    & {\color{cgpgreen}\checkmark}
    & {\color{cgpred}$\times$}
    & {\color{cgpred}$\times$}
    & {\color{cgpred}$\times$}
    & {\color{cgpred}$\times$} \\
  \rowcolor{cgplightgray}
  Changer l'ann\'ee active
    & {\color{cgpgreen}\checkmark}
    & {\color{cgpred}$\times$}
    & {\color{cgpred}$\times$}
    & {\color{cgpred}$\times$}
    & {\color{cgpred}$\times$} \\
  G\'erer les d\'el\'egations
    & {\color{cgpgreen}\checkmark}
    & {\color{cgpgreen}\checkmark}
    & {\color{cgpgreen}\checkmark}
    & {\color{cgpred}$\times$}
    & {\color{cgpred}$\times$} \\
  \rowcolor{cgplightgray}
  Soumettre demande de r\^ole
    & {\color{cgpred}$\times$}
    & {\color{cgpgreen}\checkmark}
    & {\color{cgpgreen}\checkmark}
    & {\color{cgpred}$\times$}
    & {\color{cgpred}$\times$} \\
  Traiter demandes de r\^oles
    & {\color{cgpgreen}\checkmark}
    & {\color{cgpred}$\times$}
    & {\color{cgpred}$\times$}
    & {\color{cgpred}$\times$}
    & {\color{cgpred}$\times$} \\
  \rowcolor{cgplightgray}
  Cr\'eer un signalement
    & {\color{cgpgreen}\checkmark}
    & {\color{cgpgreen}\checkmark}
    & {\color{cgpgreen}\checkmark}
    & {\color{cgpgreen}\checkmark}
    & {\color{cgpgreen}\checkmark} \\
  Traiter les signalements
    & {\color{cgpgreen}\checkmark}
    & {\color{cgporange}\textasciitilde}
    & {\color{cgporange}\textasciitilde}
    & {\color{cgpred}$\times$}
    & {\color{cgpred}$\times$} \\
  \rowcolor{cgplightgray}
  Consulter le journal d'audit
    & {\color{cgpgreen}\checkmark}
    & {\color{cgpred}$\times$}
    & {\color{cgpred}$\times$}
    & {\color{cgpred}$\times$}
    & {\color{cgpred}$\times$} \\
  Modifier son propre profil
    & {\color{cgpgreen}\checkmark}
    & {\color{cgpgreen}\checkmark}
    & {\color{cgpgreen}\checkmark}
    & {\color{cgpgreen}\checkmark}
    & {\color{cgpgreen}\checkmark} \\

\end{longtable}

\medskip
\begin{notebox}[Droits partiels (\textasciitilde)]
  La mention \og partiel \fg{} signifie que le droit existe mais est limit\'e au \textbf{p\'erim\`etre de la structure} rattach\'ee \`a l'affectation de l'utilisateur. Par exemple, un directeur de composante ne peut g\'erer les utilisateurs que dans sa composante et ses sous-structures.
\end{notebox}

\section{P\'erim\`etre et hi\'erarchie des structures}

La hi\'erarchie des structures acad\'emiques dans CGP est la suivante (du plus g\'en\'eral au plus sp\'ecifique)\,:

\medskip
\begin{center}
\begin{tikzpicture}[
  node distance=0.5cm,
  box/.style={rectangle, rounded corners=4pt, draw=cgpblue, fill=cgplightblue,
              text width=3.5cm, align=center, font=\small\bfseries, minimum height=0.8cm}
]
  \node[box] (comp) {Composante};
  \node[box, below=of comp] (dept) {D\'epartement};
  \node[box, below=of dept] (mention) {Mention};
  \node[box, below=of mention] (parcours) {Parcours};
  \node[box, below=of parcours] (niveau) {Niveau};
  \draw[->, thick, cgpblue] (comp) -- (dept);
  \draw[->, thick, cgpblue] (dept) -- (mention);
  \draw[->, thick, cgpblue] (mention) -- (parcours);
  \draw[->, thick, cgpblue] (parcours) -- (niveau);
\end{tikzpicture}
\end{center}

\medskip
Un utilisateur ayant un r\^ole dans une structure donn\'ee acc\`ede g\'en\'eralement \`a cette structure \textbf{et \`a l'ensemble de ses descendants} dans la hi\'erarchie. Les Services Centraux n'ont pas de restriction de p\'erim\`etre.

\begin{tipbox}[P\'erim\`etre de consultation versus de g\'en\'eration]
  Pour les organigrammes, la \textbf{consultation} est plus large que la \textbf{g\'en\'eration}\,: un utilisateur peut consulter des organigrammes d\'ej\`a g\'en\'er\'es pour des structures hors de son p\'erim\`etre, mais il ne peut en g\'en\'erer que pour son p\'erim\`etre propre.
\end{tipbox}
